\documentclass[10pt,a4paper]{amsart}
\usepackage[margin=19mm]{geometry}
\usepackage{amsmath,amssymb,mathtools}
\usepackage[scr=rsfs,cal=euler]{mathalfa}
\usepackage{xfrac}
\usepackage[unicode,hidelinks,bookmarks=false]{hyperref}
\newcommand{\ie}{i.e.\ }
\newcommand{\cf}{cf.\ }
\newcommand{\resp}{respectively\ }
\newcommand{\Subsection}{\subsection}
\providecommand{\nobreakdash}{\nobreak\hbox{-}}
\setlength{\emergencystretch}{2em}
\multlinegap0pt
\usepackage{bm}
\usepackage{bbm}
\usepackage{dsfont}
\usepackage{xspace}
\usepackage{cancel}
\usepackage{mathtools}
\usepackage{amsthm}
\makeatletter
\@ifundefined{proposition}{\newtheorem{proposition}{Proposition}}{}
\makeatother
\def\GL{\operatorname{GL}}
\def\ep{\epsilon}
\def\ra{\rightarrow}
\title[Ginzburg--Landau Functionals in the Large-Graph Limit]{Ginzburg--Landau Functionals in the Large-Graph Limit}
\author{Edith Zhang, James Scott, Qiang Du, Mason A. Porter}
\date{Source: arXiv:2408.00422v2 (CC BY 4.0)}
\begin{document}
\maketitle

\noindent\emph{Source and licence.} Ginzburg--Landau Functionals in the Large-Graph Limit. Version arXiv:2408.00422v2 (CC BY 4.0), distributed under the Creative Commons Attribution 4.0 International licence, as recorded in the arXiv licence metadata for this exact version and shown on the abstract page https://arxiv.org/abs/2408.00422v2. Full rights notice: licenses/arxiv-2408.00422-CC-BY-4.0.txt.\par\smallskip

\noindent\textit{Editorial scope.} Counted unit: the Proposition whose source label is \texttt{prop: community structure minimizers} in the article (source0.tex lines 1565--1568, source label \texttt{prop: community structure minimizers}), delivered together with the article's own complete original proof of it (source0.tex lines 1570--1653, one \texttt{proof} environment of 84 lines). No mathematics is written, repaired or re-derived: statement and proof are the article's own text line for line. Required context retained uncounted: eq: community graph (lines 1510--1520); eq: equivalent community fnl (lines 1538--1540); eq: communities scaling (lines 1542--1544); eq: vol constraint community (lines 1547--1549); eq: equal community vol constraint (lines 1561--1563). Every definition, display and label of those lines is retained unchanged. Omitted from this package and not redistributed: 1--1509, 1521--1537, 1541--1541, 1545--1546, 1550--1560, 1564--1564, 1569--1569, 1654--2140. Nothing inside a retained excerpt is omitted or altered: every retained body is delivered exactly as the article's lines are written, so no omission receipt of any kind accompanies this package. Citations: the retained excerpts carry no citation command at all, so no bibliography entry of the article is transcribed and no reference is redistributed. Reference adaptation: none was needed and none was made; every reference inside the delivered unit resolves inside it, so no unresolved reference or citation is produced. Printed slips kept literal: no printed word, symbol, hypothesis or number of the counted statement or of its proof was altered. Layout and class: the article's own class is \texttt{amsart} with packages including amsmath, amssymb, graphicx, amsthm, caption, cutwin, subcaption, enumerate, float, tabularx, xcolor, tikz, wrapfig, geometry; the wrapper is the lane's pinned \texttt{amsart} editorial wrapper at 10pt on A4, which restates the theorem-like environments the retained text needs and adds the article's own macro definitions that the retained text needs (\texttt{\textbackslash GL}, \texttt{\textbackslash ep}, \texttt{\textbackslash ra}), with extra wrapper packages (bm, bbm, dsfont, xspace, cancel, mathtools). The delivered numbering is the unit's own. Assets: no retained span contains a figure environment, a graphics-inclusion command, a PDF-page inclusion, a table or a TikZ picture; no image file is copied into this package, the package declares no asset dependency and no image file is redistributed with it. Licence: version arXiv:2408.00422v2 of this article is distributed under the Creative Commons Attribution 4.0 International licence, as recorded in the arXiv licence metadata for this exact version and shown on the abstract page https://arxiv.org/abs/2408.00422v2. A full rights notice is delivered at licenses/arxiv-2408.00422-CC-BY-4.0.txt. Characters: every retained excerpt is pure ASCII, so the delivered engine, XeLaTeX with its default Latin Modern text font, reports no missing character.
\begin{equation}\label{eq: community graph}
    A^{(n)}_{ij} = \begin{cases}
        1 &\text{ if }\,\, i\,, \,j \in S
        \\
        1 &\text{ if } \,\,i\,, \,j \in S^c
        \\
        0 &\text{ if } \,\,i \in S\,,\,  j\in S^c
        \\
        0 &\text{ if } \,\,i \in S^c\,,\,  j \in S \,. 
    \end{cases}
\end{equation}

\begin{equation}\label{eq: equivalent community fnl}
    E_n(v) =  \frac{c_S^4}{\ep} \sum_{i\in S} (v_i^2 - 1)^2 + \frac{c_{S^c}^4}{\ep} \sum_{i\in S^c} (v_i^2 - 1)^2 - \frac{2}{n} c_S c_{S^c} \alpha^2 - \frac{2}{n}c_S c_{S^c}(c - \alpha)^2 \,,
\end{equation}

\begin{equation}\label{eq: communities scaling}
    c_S = \sqrt{1 - \ep a} \,, \quad  c_{S^c} = \sqrt{1 - \ep(1 - a)} \,. 
\end{equation}

\begin{equation}\label{eq: vol constraint community}
    c_S \sum_{i\in S} v_i = \alpha \,, \quad c_{S^c} \sum_{i\in S^c} v_i = -\alpha \,.
\end{equation}

\begin{equation}\label{eq: equal community vol constraint}
    \frac{1}{|S|} \sum_{i\in S} v_i = \frac{\gamma}{c_S}\,, \quad  -\frac{1}{|S|} \sum_{i\in S^c} v_i = -\frac{\gamma}{c_S}\,.
\end{equation}

\begin{proposition}[Characterization of the GL minimizers for the complete community-structure graph] \label{prop: community structure minimizers}
Let $W_n$ be the complete community-structure graph, which has adjacency-matrix elements \eqref{eq: community graph}, and suppose that $|S| = |S^c|$. 
The minimizers of $\GL^{W_n}_\ep$ are functions $u$ that, on $S$, take a constant value that approaches $+1$ as $\ep \downarrow 0$ and, on $S^c$, take a constant value that approaches $-1$ as $\ep \downarrow 0$. Furthermore, the values of $\gamma$ that minimize $E_n$ {are equal to} $\pm c_S$, which approach $\pm 1$ as $\ep \ra 0$.
\end{proposition}

\begin{proof}

We separately consider the two cases $\left|{\gamma}/{c_S}\right| > 1$ and $\left|{\gamma}/{c_S}\right| \leq 1$. 

First suppose that $\left|\frac{\gamma}{c_S}\right| > 1$. 
Let $h$ be the convex hull of the double-well potential $f(s) = (s^2 - 1)^2$. Therefore, $h = 0$ on $[-1,1]$ and $h = f$ everywhere else. Define 
\begin{equation*}
    H_n(v) := \frac{2 c_S^4}{\ep} \sum_{i = 1}^n h(v_i) - \frac{4}{n} c_S^2 \alpha^2\,,
\end{equation*}
which is a lower bound of the energy $E_n$ for all $v$.
Because $h$ is convex, $(h(x) + h(y))/2 \leq h\left((x + y)/2\right)$. Therefore, if $v_i \neq v_j$ for any $i \neq j$, the function $H_n$ decreases if one replaces both $v_i$ and $v_j$ by their mean $(v_i + v_j)/2$.
Consequently, the minimizer $v'$ of $H_n$ is constant on $S$ and constant on $S^c$. However, because of the {mean-value constraint} \eqref{eq: vol constraint community}, $v'$ has different values on $S$ and $S^c$.
The minimizer that satisfies the {mean-value constraint} \eqref{eq: equal community vol constraint} is the function $v'$ with
\begin{equation}\label{eq: v'}
    v'_i = \begin{cases}
        +\frac{\gamma}{c_S} &\text{ if } \,\, i\in S 
        \\
        -\frac{\gamma}{c_S} &\text{ if } \,\, i\in S^c   \,.  
    \end{cases}
\end{equation}


Because $\left|{\gamma}/{c_S}\right| > 1$, we have that $H_n(v') = E_n(v')$. Let $\gamma' = \gamma^2$ and consider the energy
\begin{equation}\label{eq: g_1}
    E_n(v') = \frac{n c_S^4}{\ep} \left(\frac{ \gamma'}{c_S^2} - 1\right)^2 - \frac{4}{n} c_S^2 \gamma' |S|^2
\end{equation}
as a function of $\gamma'$. We denote this function by $g_1(\gamma')$.
Its derivative is
\begin{align}\label{eq: g_1 derivative}
    \frac{d}{d\gamma'}g_1(\gamma') &=  \frac{2c_S^2 n}{\ep} \left(\frac{\gamma'}{c_S^2} - 1\right) - \frac{4}{n} c_S^2 |S|^2\,.
\end{align}
Setting {the right-hand side of} \eqref{eq: g_1 derivative} to $0$ yields the critical point
\begin{equation*}
    \gamma' = c_S^2 \left(1 + \ep a \right) \,.
\end{equation*} 
Therefore, when $|S| = |S^c|$, the {critical values of $\gamma$ are}
\begin{equation*}
    \gamma \in \left \{\pm c_S \sqrt{1 + \ep/2}\right\} = \left\{\pm \sqrt{1 - (\ep/2)^2}\right\}\,. 
\end{equation*} 
The second derivative $\frac{d^2}{d\gamma'^2}g_1(\gamma') = \frac{4|S|}{\ep} > 0$ everywhere, so the critical points are minima of $g_1$. The minimizer \eqref{eq: v'} of $H_n$ is thus
\begin{equation}
    v'_i = \begin{cases}
        +\sqrt{1+\ep/2} &\text{ if }\,\, i\in S 
        \\
        -\sqrt{1+\ep/2} &\text{ if \,\,} i\in S^c   \,.  
    \end{cases}
\end{equation}
Using the change of variables \eqref{eq: communities scaling}, we obtain the GL minimizer $u$, which is defined by 
\begin{equation}
    u_i = \begin{cases}
        +\sqrt{1 - (\ep/2)^2} &\text{ if }\,\, i\in S 
        \\
        -\sqrt{1 - (\ep/2)^2} &\text{ if }\,\, i\in S^c   \,. 
    \end{cases}
\end{equation}
Because $\ep \downarrow 0$, we have that $+\sqrt{1 - (\ep/2)^2} \, \downarrow \, 1$ and $-\sqrt{1 - (\ep/2)^2} \, \uparrow \, -1$. {That is, the positive values of $u$ converge downward to $1$ and the negative values {of $u$} converge upward to $-1$.} Therefore, even when the mean value of $v$ is forced by the {mean-value constraint} \eqref{eq: equal community vol constraint} to have an absolute value larger than $1$, the optimal $u$ approaches $\pm 1$. 


Now suppose that $\left|{\gamma}/{c_S}\right| \leq 1$. We construct a candidate minimizer $\tilde{v}$ that satisfies $\gamma = \pm c_S$, and we show that the energy is larger for {$|\gamma/c_S| < 1$ than for $|\gamma/c_S| = 1$.} 
Let $\tilde{v}_i = +1$ for all nodes $i \in S$, and let $\tilde{v}_i = -1$ for all nodes $i \in S^c$. This yields $\gamma = c_S$ and an energy of
\begin{equation} \label{eq: energy bound community}
    E_n(\tilde{v}) = - \frac{4}{n} c_S^2 |S|^2 \,.
\end{equation}
Swapping the signs of $\tilde{v}$ for $S$ and $S^c$ yields $\gamma = -c_S$ and the {same energy \eqref{eq: energy bound community}}. 

Now consider $v$ with $\gamma \in (0, c_S)$, which implies that $|\gamma/c_S| < 1$. {In this case, the energy \eqref{eq: equivalent community fnl} is}
\begin{equation}\label{unsaturated energy}
    E_n(v) =  \frac{2c_S^4}{\ep} \sum_{i = 1}^n (v_i^2 - 1)^2  - \frac{4}{n} c_S^2 \gamma^2 |S|^2 \,.
\end{equation} 
The second term of \eqref{unsaturated energy} is larger than $E_n(\tilde{v})$ and the first term of \eqref{unsaturated energy} is nonnegative, so $E_n(v) \geq E_n(\tilde{v})$. 
We can use the same argument for the energy $E_n(v)$ for $\gamma \in (-c_S, 0)$. 

We conclude that ${\gamma}/{c_S} = \pm 1$ is the optimal mean value of $v$ on $S$ and $S^c$ when $\left|{\gamma}/{c_S}\right| \leq 1$. Furthermore, the candidate minimizer $\tilde{v}$ is {a} minimizer for this value of $\gamma$. 
The GL minimizer $u$ {that corresponds to this candidate minimizer $\tilde{v}$} has components
\begin{equation}
    u_i = \begin{cases}
        +\frac{1}{\sqrt{1 - \ep/2}} &\text{ if }\,\, i\in S 
        \\
        -\frac{1}{\sqrt{1 - \ep/2}} &\text{ if } \,\, i\in S^c \,. 
    \end{cases}
\end{equation}
{As $\ep \downarrow 0$, the values of the {GL} minimizer $u$ approach $\pm 1$.}

\end{proof}

\end{document}
